\section{Problem Formulation}
\label{sec:problem-formulation}

The main quantities are client subscriptions $\mathcal{K}_u$, active sockets
$\mathcal{C}_u$, per-subscription backlog $b_{u,k}$, and the corresponding
connection loads $(N_{u,c},G_{u,c})$. Pod-wide quantities $(X_i,B_i)$ include
all clients served by Pod $i$. Appendix~\suppref{app:notation} collects the full
symbol reference; definitions also appear at their first use below.

We model the system at discrete control epochs $t=0,1,\ldots,T$. The interval
between adjacent epochs is $\Delta$ seconds: the next epoch is $t+1$, while
its physical time is $(t+1)\Delta$. Live-event timestamps and durations are
reported separately in seconds. Let
$\mathcal{S}=\{1,\ldots,n\}$ be the set of subscriber pods and
$\mathcal{U}=\{1,\ldots,m\}$ be the set of clients. Client $u$ maintains a
subscription set $\mathcal{K}_u(t)$ and a set of persistent transport
connections $\mathcal{C}_u(t)$, with
$1\leq |\mathcal{C}_u(t)|\leq C_{\max}$.
Indices $i,u,c,k$ refer to Pods, clients, connections, and subscriptions,
respectively. Calligraphic symbols denote sets; their cardinalities are
counts. In the single-client scaling rules, $C(t)=|\mathcal{C}_u(t)|$ and
$N_{\mathrm{sub}}(t)=|\mathcal{K}_u(t)|$ are count shorthands with the fixed client
index suppressed. They are never used as sets. Figure labels and algorithm
steps may omit $(t)$ when all quantities refer to the same decision epoch.

\subsection{Per-Subscription Queues}

Each subscription has a queue that stores messages waiting to be delivered and
processed. Let $Q_{u,k}(t)$ be the queue of subscription
$k\in\mathcal{K}_u(t)$, and let
\begin{equation}
	b_{u,k}(t)=|Q_{u,k}(t)|
\end{equation}
be the backlog for that subscription. This is the basic backlog measure used
by the client. At the pod level, backlog is a
server-owned aggregate. Let $\mathcal{R}_i(t)$ denote the set of subscription
queues currently served by pod $i$, and let $\tilde{b}_{i,r}(t)$ be the backlog
of server-side subscription queue $r\in\mathcal{R}_i(t)$. If queue $r$
corresponds to client $u$ and subscription $k$, then $\tilde{b}_{i,r}(t)$ is
the server-side counterpart of $b_{u,k}(t)$; however, $\mathcal{R}_i(t)$ also
contains queues belonging to other clients and is maintained by the pod, not by
any single client.

Let $z_{u,k,c}(t)\in\{0,1\}$ indicate whether subscription $k$ of client $u$ is
assigned to connection $c\in\mathcal{C}_u(t)$. Feasible assignment requires
\begin{equation}
	\sum_{c\in\mathcal{C}_u(t)} z_{u,k,c}(t)=1,
	\qquad \forall u,\; k\in\mathcal{K}_u(t).
	\label{eq:subscription_assignment}
\end{equation}
Let $i_{u,c}(t)\in\mathcal{S}$ denote the pod index reached by connection
$c$. This endpoint mapping differs from the Pod name $p_i$ and from the
selected DQL action $a_u(t)$.
The number of subscriptions and cumulative subscription backlog on connection
$c$ are
\begin{align}
	N_{u,c}(t) &= \sum_{k\in\mathcal{K}_u(t)} z_{u,k,c}(t),\\
	G_{u,c}(t) &= \sum_{k\in\mathcal{K}_u(t)}
	b_{u,k}(t)z_{u,k,c}(t).
	\label{eq:connection_backlog}
\end{align}
The first term captures subscription-management pressure on the connection; the
second captures queue pressure induced by the subscriptions mapped to that
connection.

\subsection{Pod Statistics and Partial Observability}

Each pod exposes a statistics tuple through active connections and probe
connections:
\begin{equation}
	y_i(t)=\bigl(p_i,\; X_i(t),\; B_i(t),\; H_i\bigr),
	\label{eq:pod_stats}
\end{equation}
Here $p_i$ is the pod name, and $X_i(t)$ is its subscription count.
The value $B_i(t)$ is the cumulative backlog across all subscription queues
served by pod $i$. The value $H_i$ is the configured maximum processing threshold. Using the
server-side queue set $\mathcal{R}_i(t)$, the pod-level subscription count and
cumulative pod backlog are
\begin{equation}
	X_i(t)=|\mathcal{R}_i(t)|,\qquad
	B_i(t)=
	\sum_{r\in\mathcal{R}_i(t)}
	\tilde{b}_{i,r}(t).
	\label{eq:pod_backlog}
\end{equation}
Thus $B_i(t)$ is not defined by summing over the connection set
$\mathcal{C}_u(t)$ of a particular client. A client's connections determine
which pod statistics it can observe and which subscriptions it can move, but
the pod backlog itself is the cumulative backlog over all subscription queues
served by that pod. These pod-level values include traffic from many clients
and are not directly computable by a single client unless reported by the pod.

At each statistics interval $\Delta$ (10 seconds in the target deployment),
every active connection receives the tuple of its current pod. In addition,
client $u$ maintains one authenticated probe connection whose endpoint is
selected by the load balancer and receives its tuple. The client
retains the probe socket and its endpoint identity until promotion, expiry, or
recycling; it does not resample a different pod merely by reading the next
observation. Therefore the client observes only a subset
$\mathcal{O}_u(t)\subseteq\mathcal{S}$ of pod states:
\begin{equation}
	\mathcal{O}_u(t)=
	\{i_{u,c}(t):c\in\mathcal{C}_u(t)\}
	\cup
	\{i^{\mathrm{probe}}_u(t)\}.
\end{equation}
The client stores the newest tuple for each observed pod in a local cache. This
makes the placement problem partially observable: global cluster variance is
available to an offline evaluator or simulator, but not to an individual client
controller. In particular, an admissible migration target must be represented
by a currently retained probe socket; the client cannot replace the sampled
candidate with an unobserved globally optimal pod.

\begin{assumption}[Available deployment signals]
For every pod $i$ observed by client $u$, where $i\in\mathcal{O}_u(t)$, the
protocol supplies $(p_i,X_i(t),B_i(t),H_i)$ in a statistics frame, header, or
control message. Client
$u$ can also measure its own per-subscription backlogs $b_{u,k}(t)$, delivered
message count, processed message count, acknowledged message count, active
connection count, and migration count over each interval $\Delta$.
\end{assumption}

\begin{assumption}[Transport abstraction]
The underlying transport is required only to maintain long-lived client--pod
connections, multiplex one or more subscriptions over a connection, support
subscription migration or reassignment, and expose periodic statistics through
headers, frames, or control messages. STOMP-over-TCP, custom TCP streams, and
WebSocket-based messaging are concrete examples of this abstraction.
\end{assumption}

\begin{remark}
The pod threshold $H_i$ is not treated as an actual drain rate. The true drain
rate depends on the queue subsystem, pod scheduling, network transport,
connection behavior, client processing capacity, and acknowledgement rate.
Hence $B_i(t)/(H_i\Delta+\varepsilon_0)$ is used only as a dimensionless
pod-side pressure ratio. Here $B_i$ counts messages, $H_i$ has units of
events per second, and $\varepsilon_0$ has units of messages. The ratio is
not an estimate of the actual drain time.
\end{remark}

\subsection{Local Imbalance Measures}

For client $u$, define the mean subscription count and mean backlog load across
its own connections:
\begin{align}
	\bar{N}_u(t)
	&= \frac{1}{|\mathcal{C}_u(t)|}
	\sum_{c\in\mathcal{C}_u(t)} N_{u,c}(t),\\
	\bar{G}_u(t)
	&= \frac{1}{|\mathcal{C}_u(t)|}
	\sum_{c\in\mathcal{C}_u(t)} G_{u,c}(t).
\end{align}
The corresponding variances are
\begin{align}
	V^N_u(t)
	&= \frac{1}{|\mathcal{C}_u(t)|}
	\sum_{c\in\mathcal{C}_u(t)}
	\bigl(N_{u,c}(t)-\bar{N}_u(t)\bigr)^2,\\
	V^G_u(t)
	&= \frac{1}{|\mathcal{C}_u(t)|}
	\sum_{c\in\mathcal{C}_u(t)}
	\bigl(G_{u,c}(t)-\bar{G}_u(t)\bigr)^2.
\end{align}
The first variance captures subscription-management skew, while the second
captures backlog skew. The local rebalancing objective is therefore
\begin{equation}
	V_u(t)=
	\eta_N\widehat{V}^N_u(t)
	+\eta_G\widehat{V}^G_u(t),
	\qquad
	\eta_N+\eta_G=1,
	\label{eq:local_connection_imbalance}
\end{equation}
where hats denote normalized quantities. This formulation is necessary because
two connections can have equal backlog but very different subscription counts,
or equal subscription counts but very different backlog.
The normalized terms place the two variances on a common scale before they are
combined. In the evaluation setup, we use positive reference values
$V^N_{\mathrm{ref}}$ and $V^G_{\mathrm{ref}}$ and compute
\begin{align*}
	\widehat{V}^N_u(t)
	&=\min\left\{\frac{V^N_u(t)}{V^N_{\mathrm{ref}}},1\right\},\\
	\widehat{V}^G_u(t)
	&=\min\left\{\frac{V^G_u(t)}{V^G_{\mathrm{ref}}},1\right\}.
\end{align*}
Thus $V_u(t)$ is a dimensionless connection-imbalance score in which the
subscription-count component and the backlog-load component are comparable.
The weights $\eta_N,\eta_G\geq0$ determine how much the controller prioritizes
subscription-management fairness versus backlog-load fairness. For example,
larger $\eta_N$ makes the controller more sensitive to connections carrying
too many subscriptions, while larger $\eta_G$ makes it more sensitive to
connections accumulating more queued messages. The constraint
$\eta_N+\eta_G=1$ makes $V_u(t)$ a convex combination rather than an
unbounded sum.

\subsection{Pod Pressure and Deployable Reward}

For an observed pod $i\in\mathcal{O}_u(t)$, define normalized pressure terms
\begin{align}
	P^X_i(t) &= \min\left\{\frac{X_i(t)}{X_{\mathrm{safe}}},1\right\},\\
	P^B_i(t) &= \min\left\{\frac{B_i(t)}{B_{\mathrm{ref}}},1\right\},\\
	P^H_i(t) &= \min\left\{\frac{B_i(t)}{H_i\Delta+\varepsilon_0},1\right\},
	\qquad \varepsilon_0>0.
\end{align}
Here $X_{\mathrm{safe}}$ is the subscription count at which pod reliability begins to
degrade. The value $X_{\mathrm{hard}}$ is the hard reference count at which we treat the
pod as severely overloaded by subscription-management work.
The value $B_{\mathrm{ref}}$ is the backlog reference used for normalization. The condition
$X_{\mathrm{hard}}>X_{\mathrm{safe}}$ is required. Since high subscription count has been
observed to make connections unreliable, we also define a nonlinear
reliability-risk penalty
\begin{equation}
	P^R_i(t)=
	\left[
	\frac{X_i(t)-X_{\mathrm{safe}}}
	{X_{\mathrm{hard}}-X_{\mathrm{safe}}}
	\right]_+^2,
	\qquad [x]_+=\max\{x,0\}.
\end{equation}
The pod pressure score is
\begin{align}
	P_i(t)=
	&\omega_X P^X_i(t)
	+\omega_B P^B_i(t) \nonumber\\
	&+\omega_H P^H_i(t)
	+\omega_R P^R_i(t),
	\label{eq:pod_pressure}\\
	&\sum_{q\in\{X,B,H,R\}}\omega_q=1,\qquad \omega_q\geq 0. \nonumber
\end{align}
Here $P^X_i(t)$ measures normalized subscription-count pressure.
The term $P^B_i(t)$ measures normalized cumulative-backlog pressure.
The term $P^H_i(t)$ measures backlog pressure relative to the processing threshold.
Finally, $P^R_i(t)$ measures nonlinear reliability risk and becomes active only
when $X_i(t)>X_{\mathrm{safe}}$. The coefficients
$\omega_X,\omega_B,\omega_H,\omega_R$ specify how much importance the policy
assigns to each pressure component. For example, if subscription-heavy pods are
known to become unreliable before backlog becomes large, then $\omega_X$ and
$\omega_R$ receive greater weight than $\omega_B$ and $\omega_H$. If queue
accumulation is the dominant failure mode, $\omega_B$ and $\omega_H$
receive greater weight. This score keeps subscription count and backlog
count as distinct load dimensions.

Client $u$ also computes local backlog and backlog-growth pressure:
\begin{align}
	B_u(t) &= \sum_{k\in\mathcal{K}_u(t)} b_{u,k}(t),\\
	\mathrm{Grow}_u(t) &= \bigl[B_u(t)-B_u(t-1)\bigr]_+,\qquad t\geq1.
\end{align}
At initialization, $\mathrm{Grow}_u(0)=0$. Let $T_u(t)$ be the effective client-side
throughput in events per second. The rates $T^{\mathrm{del}}_u(t)$, $T^{\mathrm{proc}}_u(t)$,
and $T^{\mathrm{ack}}_u(t)$ are the delivered, processed, and acknowledged counts in
the measured interval, each divided by its duration $\Delta$:
\begin{equation}
	T_u(t)=
	\min\{T^{\mathrm{del}}_u(t),T^{\mathrm{proc}}_u(t),T^{\mathrm{ack}}_u(t)\}.
\end{equation}
For an action that selects or keeps pod $i$, the deployable reward is
\begin{align}
	R_u(t)=
	&-\lambda_1\widehat{B}_u(t)
	-\lambda_2\widehat{\mathrm{Grow}}_u(t)
	-\lambda_3 P_i(t)
	-\lambda_4 V_u(t) \nonumber\\
	&-\lambda_5 M_u(t)
	-\lambda_6\widehat{C}_u(t)
	+\lambda_7\widehat{T}_u(t),
	\label{eq:deployable_reward}
\end{align}
where $M_u(t)$ is the migration indicator or normalized migration count,
$\widehat{C}_u(t)=|\mathcal{C}_u(t)|/C_{\max}$, and
$\lambda_\ell\geq0$ with $\sum_{\ell=1}^{7}\lambda_\ell=1$.
The hats on backlog, backlog growth, and throughput denote their normalized
versions under the stated experiment's scaling convention, not different
physical quantities. This reward is computable by the client from
local measurements and observed pod tuples.

Detailed reference-scale calibration, the global evaluation objective, and
step-by-step proofs of the subscription-density and backlog-variance
properties are provided in Section~\suppref{app:extended-mathematics}. The full notation reference is in
Section~\suppref{app:notation}.
